\documentclass[11pt]{article}

\usepackage{amsmath,amssymb,amsthm,bm,mathtools}
\usepackage{geometry}
\geometry{margin=1in}
\usepackage{setspace}
\usepackage{enumitem}
\usepackage{booktabs}
\usepackage{hyperref}
\usepackage{algorithm}
\usepackage{algpseudocode}
\usepackage{float}

\onehalfspacing

\newcommand{\R}{\mathbb R}
\newcommand{\ind}{\mathbb I}
\newcommand{\norm}[1]{\left\lVert #1 \right\rVert}
\newcommand{\trans}{^{\mathsf T}}
\DeclareMathOperator{\logit}{logit}
\DeclareMathOperator{\Cat}{Categorical}
\DeclareMathOperator{\Dir}{Dirichlet}
\DeclareMathOperator{\IG}{IG}
\DeclareMathOperator{\IW}{IW}
\DeclareMathOperator{\tr}{tr}

\title{Joint Multilayered LSIRM with Hierarchical Latent-Position Gaussian Mixture Clustering\\
Cluster-Aligned Auxiliary Positions and Decoupled Anchor Covariance}
\date{Model Description}

\begin{document}
\maketitle

\section{Overview}

This document defines a joint multilayered LSIRM in which item clustering is
performed on \emph{cluster-aligned auxiliary latent positions} that sit
between the LSIRM item positions and the Gaussian-mixture cluster centres.
The model is a hierarchical refinement of the latent-position mixture prior
of De~Carolis, Kang, and Jeon (2025): instead of placing a finite Gaussian
mixture directly on each item's LSIRM position $b_j^{(\ell)}$, we introduce
a per-item auxiliary position $r_q\in\R^d$ for each global item $q$ and
specify
\[
    b_j^{(\ell(q))}=z_q\mid r_q,\Sigma_b\sim N_d(r_q,\Sigma_b),
    \qquad
    r_q\mid c_q=l,\mu_l,\Sigma_l\sim N_d(\mu_l,\Sigma_l),
\]
with the same NIW prior on $(\mu_l,\Sigma_l)$ as before but a separate global
within-anchor covariance $\Sigma_b$ that does \emph{not} appear in the cluster
component definition.

\paragraph{Motivation.}
In the unhierarchical version, the cluster covariance $\Sigma_{c_q}$ plays
two roles simultaneously: it defines cluster shape (via the mixture component
$N(\mu_l,\Sigma_l)$) and it sets the prior precision pulling $b_j^{(\ell)}$
toward the cluster centre in the Metropolis--Hastings update of the item
position.  Tight clusters give small $\Sigma_l$, which compresses
within-cluster pairwise distances among item positions, distorting the LSIRM
geometry that downstream analyses report.  Splitting the role into a
cluster-shape covariance $\Sigma_l$ (on $r_q$) and a global anchor
covariance $\Sigma_b$ (between $z_q$ and $r_q$) leaves cluster shape
heterogeneity intact while pooling anchor strength across all $P$ items, so
that anchor pull no longer collapses with cluster tightness.

\paragraph{Marginal cluster-conditional distribution.}
Integrating $r_q$ out of the hierarchy,
\begin{equation}
    z_q\mid c_q=l,\mu_l,\Sigma_l,\Sigma_b\sim N_d(\mu_l,\;\Sigma_l+\Sigma_b),
    \label{eq:z-marginal}
\end{equation}
which makes explicit that cluster shape heterogeneity is carried by the
per-cluster $\Sigma_l$, while $\Sigma_b$ adds a uniform additional dispersion
common to all clusters.  Cluster membership $c_q$ is an item-level label
that applies equally to $r_q$ and $z_q$; visualizations and posterior
similarity matrices over items are unchanged in interpretation.

\paragraph{Modelling restriction.}
Cluster labels affect $r_q$, and through it $z_q=b_{j(q)}^{(\ell(q))}$, but
no cluster-specific prior is placed on respondent ability, item difficulty,
or layer-specific nuisance parameters.

\section{Indexing, dimensions, and observed data}

\begin{center}
\begin{tabular}{lll}
\toprule
Symbol & Meaning & Dimension/range \\
\midrule
$\ell$ & layer index & $\ell=1,\dots,L$ \\
$i$ & respondent/person index & $i=1,\dots,n_\ell$ \\
$j$ & local item index within layer $\ell$ & $j=1,\dots,P_\ell$ \\
$q$ & global item index across all layers & $q=1,\dots,P$, $P=\sum_{\ell=1}^L P_\ell$ \\
$\ell(q)$ & layer containing global item $q$ & $\{1,\dots,L\}$ \\
$j(q)$ & local item index of global item $q$ & $\{1,\dots,P_{\ell(q)}\}$ \\
$d$ & latent space dimension & $d=2$ in the current implementation \\
$K^\star$ & truncation level / maximum number of mixture components & positive integer \\
$K_+$ & number of occupied mixture components & $K_+=|\{l:n_l>0\}|$ \\
$z_q$ & global item latent position (LSIRM-side) & $z_q=b_{j(q)}^{(\ell(q))}\in\R^d$ \\
$r_q$ & cluster-aligned auxiliary latent position & $r_q\in\R^d$ \\
$c_q$ & cluster label of item $q$ & $c_q\in\{1,\dots,K^\star\}$ \\
$\mu_l,\Sigma_l$ & cluster $l$ mean and shape covariance & $\mu_l\in\R^d$, $\Sigma_l\in\R^{d\times d}$ pos.\ def. \\
$\Sigma_b$ & global anchor covariance ($z_q$ around $r_q$) & $\Sigma_b\in\R^{d\times d}$ pos.\ def. \\
$S_0,S_b$ & NIW / IW scale matrices for $\Sigma_l,\Sigma_b$ & $\R^{d\times d}$ pos.\ def. \\
\bottomrule
\end{tabular}
\end{center}

\paragraph{Common-respondent assumption.}
As in the unhierarchical version, $n_\ell=n$ for all $\ell$, and the
respondent latent position is shared across layers,
$a_i^{(\ell)}=a_i\in\R^d$.  The layered notation $a_i^{(\ell)}$ is kept
for visual symmetry with $b_j^{(\ell)}$ but should be read as the single
$a_i$.  The Metropolis--Hastings update of $a_i$ in
\eqref{eq:a-mh-ratio} aggregates likelihood contributions across all layers
that contain respondent $i$, with a single fixed unit-variance prior
\eqref{eq:a-prior}.

The observed data are denoted by
$Y^{(\ell)}=\{y_{ij}^{(\ell)}:(i,j)\in\mathcal O_\ell\}$, where
$\mathcal O_\ell$ is the set of observed respondent--item pairs in layer
$\ell$.  If there is no missingness, then
$\mathcal O_\ell=\{1,\dots,n\}\times\{1,\dots,P_\ell\}$.

\section{LSIRM likelihood}

The LSIRM likelihood is unchanged from the unhierarchical model.  For
respondent $i$ and item $j$ in layer $\ell$, the linear predictor is
\begin{equation}
    \eta_{ij}^{(\ell)}
    =
    \alpha_i^{(\ell)}-\beta_j^{(\ell)}-\gamma_\ell\norm{a_i-b_j^{(\ell)}}_2,
    \qquad \ell\notin\mathcal L_{\mathrm{ord}},
    \label{eq:lsirm-linear-predictor}
\end{equation}
with the convention that for ordinal layers $\beta_j^{(\ell)}$ does not appear
in the linear predictor:
\begin{equation}
    \eta_{ij}^{(\ell)}
    =
    \alpha_i^{(\ell)}-\gamma_\ell\norm{a_i-b_j^{(\ell)}}_2,
    \qquad \ell\in\mathcal L_{\mathrm{ord}}.
    \label{eq:lsirm-linear-predictor-ord}
\end{equation}
The per-layer observation models (binary, continuous, count, ordinal LSGRM)
are exactly as specified in the original document and are summarized
compactly as
\begin{equation}
    y_{ij}^{(\ell)}\mid \eta_{ij}^{(\ell)},\xi_\ell
    \sim f_\ell\bigl\{y_{ij}^{(\ell)}\mid \eta_{ij}^{(\ell)},\xi_\ell\bigr\},
    \qquad (i,j)\in\mathcal O_\ell,
    \label{eq:generic-layer-likelihood}
\end{equation}
with the layer-specific nuisance parameter $\xi_\ell$ collecting
$(\sigma_0^2,\lambda_{ij}^{(\ell)})$ in continuous layers,
$\{\kappa_j\}_{j=1}^{P_\ell}$ in count layers, the threshold vector
$(\beta_{j,1}^{(\ell)},\dots,\beta_{j,K_\ell-1}^{(\ell)})$ in ordinal layers,
and $\gamma_\ell$ in every layer.  The full LSIRM log-likelihood is
\begin{equation}
    \ell_{\mathrm{LSIRM}}
    =
    \sum_{\ell=1}^L
    \sum_{(i,j)\in\mathcal O_\ell}
    \log f_\ell\bigl\{y_{ij}^{(\ell)}\mid \eta_{ij}^{(\ell)},\xi_\ell\bigr\}.
    \label{eq:lsirm-loglik}
\end{equation}
The crucial dependence on the clustering machinery enters through
$b_j^{(\ell)}=z_q$ in the distance term of
\eqref{eq:lsirm-linear-predictor}--\eqref{eq:lsirm-linear-predictor-ord}.

\section{Hierarchical latent-position mixture prior for items}

\subsection{Two-level item position hierarchy}

For each global item index $q=1,\dots,P$, we introduce two latent positions
in $\R^d$:
\begin{itemize}[leftmargin=1.5em]
    \item $z_q=b_{j(q)}^{(\ell(q))}$, the LSIRM item position appearing in
    \eqref{eq:lsirm-linear-predictor}--\eqref{eq:lsirm-linear-predictor-ord}.
    \item $r_q$, an auxiliary cluster-aligned position that mediates the
    relation between $z_q$ and the mixture component centred at $\mu_{c_q}$.
\end{itemize}

The hierarchy is
\begin{align}
    z_q\mid r_q,\Sigma_b
    &\sim N_d(r_q,\Sigma_b),
    \label{eq:z-given-r}\\
    r_q\mid c_q=l,\mu_l,\Sigma_l
    &\sim N_d(\mu_l,\Sigma_l),
    \label{eq:r-given-cluster}\\
    c_q\mid \rho
    &\sim \Cat(\rho_1,\dots,\rho_{K^\star}),\\
    \rho
    &\sim \Dir(e_0,\dots,e_0).
    \label{eq:rho-prior}
\end{align}

\subsection{NIW prior on cluster parameters and IW prior on anchor covariance}

The mixture component parameters carry the same NIW prior as in the
unhierarchical model:
\begin{align}
    \Sigma_l
    &\sim \IW(\nu_0,S_0),
    \label{eq:sigma-niw-prior}\\
    \mu_l\mid \Sigma_l
    &\sim N_d\!\left(m_0,\frac{1}{\kappa_0}\Sigma_l\right),
    \qquad l=1,\dots,K^\star,
    \label{eq:mu-niw-prior}
\end{align}
with $m_0\in\R^d$, $\kappa_0>0$, $\nu_0>d-1$, and $S_0\in\R^{d\times d}$
positive definite.

The new global anchor covariance has its own IW prior, independent of the
cluster NIW:
\begin{equation}
    \Sigma_b\sim\IW(\nu_b,S_b),
    \qquad \nu_b>d-1,\;S_b\in\R^{d\times d}\;\text{positive definite}.
    \label{eq:sigma-b-prior}
\end{equation}

\paragraph{Wishart hyperpriors on $S_0$ and $S_b$.}
Both scale matrices may be learned from data via Wishart hyperpriors:
\begin{align}
    S_0&\sim W(\nu_{S_0},\Lambda_{S_0}),
    \qquad \nu_{S_0}>d-1,\;\Lambda_{S_0}\in\R^{d\times d}\;\text{positive definite},
    \label{eq:S0-hyperprior}\\
    S_b&\sim W(\nu_{S_b},\Lambda_{S_b}),
    \qquad \nu_{S_b}>d-1,\;\Lambda_{S_b}\in\R^{d\times d}\;\text{positive definite}.
    \label{eq:Sb-hyperprior}
\end{align}
Both are conjugate for their respective IW scale matrices and yield closed
form Wishart full conditionals.  Defaults are weakly informative and
anchored at user-provided initial values:
$\nu_{S_0}=d+2$, $\Lambda_{S_0}=S_0^{(0)}/\nu_{S_0}$,
$\nu_{S_b}=d+2$, $\Lambda_{S_b}=S_b^{(0)}/\nu_{S_b}$, yielding
$E[S_0]=S_0^{(0)}$ and $E[S_b]=S_b^{(0)}$.

\subsection{Identifiability of $\Sigma_l$ versus $\Sigma_b$}
\label{sec:identifiability}

Marginalizing $r_q$ from
\eqref{eq:z-given-r}--\eqref{eq:r-given-cluster} gives
\eqref{eq:z-marginal}: $z_q\mid c_q=l\sim N_d(\mu_l,\Sigma_l+\Sigma_b)$.
Within a single cluster, only the sum $\Sigma_l+\Sigma_b$ is identifiable
from the LSIRM likelihood through $z$ alone.  Two pieces of information
separate $\Sigma_l$ from $\Sigma_b$ in the joint posterior:
\begin{enumerate}[leftmargin=1.5em]
    \item $\Sigma_b$ is global and contributes
    $\sum_{q=1}^P\log N_d(z_q;r_q,\Sigma_b)$ to the joint, accumulating
    information from all $P$ items.
    \item $\Sigma_l$ is per-cluster and is informed only by the items
    currently assigned to cluster $l$.
\end{enumerate}
Conditional on $r$, both are identified, and the priors
\eqref{eq:sigma-niw-prior}, \eqref{eq:sigma-b-prior} regularize the
decomposition.

\paragraph{Identifiability in small clusters.}
The decomposition relies on multi-cluster pooling through the global
$\Sigma_b$.  When a cluster has very few items ($n_l\leq 2$, including
the $K_+=1$ degenerate case), $\Sigma_l$ is informed almost entirely by
its NIW prior, and $\Sigma_b$ can absorb most of the within-cluster
variance without penalty from the likelihood.  The diagnostic
$\tr(\Sigma_b)/\min_l\tr(\Sigma_l)$ should be monitored: a sustained
divergence ($\gtrsim 10$) signals that $\Sigma_b$ is absorbing cluster
shape rather than acting as a global LSIRM-side anchor noise term, and
either the prior on $\Sigma_b$ should be tightened or the prior on
$\Sigma_l$ for small clusters should be relaxed.

\paragraph{Recommended default: isotropic $\Sigma_b$.}
For maximum identifiability and ease of interpretation, we recommend
restricting $\Sigma_b=\sigma_b^2 I_d$ as a default, deferring anisotropic
$\Sigma_b$ to applications where it is empirically warranted.  The
restriction lets the cluster-side $\Sigma_l$ carry all anisotropy and
reduces $\Sigma_b$ to a scalar ``LSIRM-side noise level''.  The
corresponding prior is
\begin{equation}
    \sigma_b^2\sim\IG(a_b,b_b),
    \qquad a_b>0,\;b_b>0,
    \label{eq:sigma-b-scalar-prior}
\end{equation}
with full conditional in \eqref{eq:sigma-b-scalar-update}.  All other
clustering machinery (NIW on $(\mu_l,\Sigma_l)$, Student-$t$ predictive
collapsed allocation, split--merge moves) is unaffected by this choice.
The full IW form \eqref{eq:sigma-b-prior} should be used when the data
provide strong evidence for anisotropic LSIRM-side noise.

\section{Priors for non-clustering LSIRM parameters}

These priors are unchanged from the unhierarchical model.

\paragraph{Respondent latent position (fixed unit variance).}
\begin{equation}
    a_i \sim N_d(0,I_d),\qquad i=1,\dots,n.
    \label{eq:a-prior}
\end{equation}

\paragraph{Respondent effect (layer-specific variance, sampled).}
\begin{equation}
    \alpha_i^{(\ell)}\mid \sigma_{\alpha,\ell}^2
    \sim N(0,\sigma_{\alpha,\ell}^2),
    \qquad
    \sigma_{\alpha,\ell}^2 \sim \IG(A_{\alpha,\ell},B_{\alpha,\ell}).
    \label{eq:alpha-prior}
\end{equation}

\paragraph{Item difficulty (fixed-scale prior).}
For non-ordinal layers,
$\beta_j^{(\ell)}\sim N(0,\sigma_{\beta,\ell}^2)$ with $\sigma_{\beta,\ell}^2$
fixed.  For ordinal layers, each threshold satisfies
$\beta_{j,k}^{(\ell)}\sim N(0,\sigma_{\beta,\ell}^2)$ subject to
$\beta_{j,1}^{(\ell)}<\dots<\beta_{j,K_\ell-1}^{(\ell)}$.

The item LSIRM position $b_j^{(\ell)}=z_q$ does not receive an
independent Gaussian prior; its prior is the hierarchy in
\eqref{eq:z-given-r}--\eqref{eq:r-given-cluster} together with
\eqref{eq:sigma-niw-prior}--\eqref{eq:sigma-b-prior}.

The remaining nuisance parameters --- $\gamma_\ell$, $\sigma_0^2$,
$\{\kappa_j\}$ --- are equipped with the priors specified in the
unhierarchical model.

\section{Rotation and reflection invariance}
\label{sec:layer-comparability}

Under the common-respondent assumption, the entire configuration
$(a,b^{(1)},\dots,b^{(L)},r)$ shares one coordinate system.  The only
identifiability ambiguity in the latent space is a single global
orthogonal transformation $Q\in O(d)$ together with a single global
translation $t\in\R^d$:
\begin{equation}
    a_i\mapsto Qa_i+t,
    \quad
    b_j^{(\ell)}\mapsto Qb_j^{(\ell)}+t,
    \quad
    r_q\mapsto Qr_q+t,
    \quad
    \mu_l\mapsto Q\mu_l+t,
    \quad
    \Sigma_l\mapsto Q\Sigma_l Q\trans,
    \quad
    \Sigma_b\mapsto Q\Sigma_b Q\trans.
    \label{eq:rigid-motion}
\end{equation}
The LSIRM linear predictor depends on positions only through
$\norm{a_i-b_j^{(\ell)}}_2$, which is invariant.  With $m_0=0$,
$S_0=s_0 I_d$, and $S_b=s_b I_d$ (or $\Sigma_b=\sigma_b^2 I_d$), the priors
on $(\mu_l,\Sigma_l)$ and $\Sigma_b$ are also invariant in distribution
under $Q\in O(d)$.  Therefore the joint posterior on the clustering
targets---$c$, $K_+$, the cluster sizes $n_l$, and the posterior
similarity matrix $C_{qr}=P(c_q=c_r\mid Y)$---is invariant under
\eqref{eq:rigid-motion}.

\section{Full joint posterior}

Let
\begin{align*}
    \Theta_{\mathrm{LSIRM}}
    &=\{a_i,\alpha_i^{(\ell)},b_j^{(\ell)},\beta_j^{(\ell)},\xi_\ell,
    \sigma_{\alpha,\ell}^2:i,\ell\},\\
    \Theta_{\mathrm{mix}}
    &=\{r_q,c_q,\rho,\mu_l,\Sigma_l:q=1,\dots,P,\;l=1,\dots,K^\star\}\cup\{\Sigma_b\}.
\end{align*}
The full posterior is proportional to
\begin{align}
&p(\Theta_{\mathrm{LSIRM}},\Theta_{\mathrm{mix}}\mid Y)
\notag\\
&\quad \propto
\prod_{\ell=1}^L
\prod_{(i,j)\in\mathcal O_\ell}
    f_\ell\{y_{ij}^{(\ell)}\mid \eta_{ij}^{(\ell)},\xi_\ell\}
\notag\\
&\qquad\times
\prod_{q=1}^P
    N_d(z_q;r_q,\Sigma_b)\,
    N_d(r_q;\mu_{c_q},\Sigma_{c_q})\,\rho_{c_q}
\notag\\
&\qquad\times
\prod_{l=1}^{K^\star}p(\mu_l\mid\Sigma_l)p(\Sigma_l)
\;p(\rho)\;p(\Sigma_b)
\notag\\
&\qquad\times
\prod_{i=1}^{n}N_d(a_i;0,I_d)
\prod_{\ell=1}^L
\left[
\prod_{i=1}^{n}p(\alpha_i^{(\ell)}\mid\sigma_{\alpha,\ell}^2)
\;p_\beta^{(\ell)}\bigl(\beta^{(\ell)}\bigr)
\right]
\notag\\
&\qquad\times
\prod_{\ell=1}^L p(\sigma_{\alpha,\ell}^2)p(\xi_\ell),
    \label{eq:full-posterior}
\end{align}
where $z_q=b_{j(q)}^{(\ell(q))}$ and the layer-specific item-difficulty prior
factor $p_\beta^{(\ell)}$ is as in the unhierarchical model.  The
hyperpriors on $S_0,S_b$ contribute factors $p(S_0),p(S_b)$ when used.

\section{NIW posterior quantities (on $r$)}

For a generic set of global item indices $A\subseteq\{1,\dots,P\}$, define
\begin{align}
    n_A &= |A|,\\
    \bar r_A &= \frac{1}{n_A}\sum_{q\in A}r_q,
    \qquad n_A>0,\\
    Q_A &= \sum_{q\in A}(r_q-\bar r_A)(r_q-\bar r_A)\trans,
    \qquad n_A>0.
\end{align}
For $n_A=0$, use prior quantities directly and set $Q_A=0$.

The NIW posterior parameters for the items in $A$ are
\begin{align}
    \kappa_A &= \kappa_0+n_A,\\
    \nu_A &= \nu_0+n_A,\\
    m_A &=
    \begin{cases}
    \dfrac{\kappa_0m_0+n_A\bar r_A}{\kappa_0+n_A}, & n_A>0,\\[1em]
    m_0, & n_A=0,
    \end{cases}\\
    S_A &=
    \begin{cases}
    S_0+Q_A+\dfrac{\kappa_0 n_A}{\kappa_0+n_A}(\bar r_A-m_0)(\bar r_A-m_0)\trans,
    & n_A>0,\\[1em]
    S_0, & n_A=0.
    \end{cases}
    \label{eq:niw-posterior-params}
\end{align}

For component $l$, let
\[
    \mathcal J_l=\{q:c_q=l\},
    \qquad n_l=|\mathcal J_l|,
\]
and set $(\kappa_l,\nu_l,m_l,S_l)=(\kappa_{\mathcal J_l},\nu_{\mathcal J_l},m_{\mathcal J_l},S_{\mathcal J_l})$.
For leave-one-out calculations, $\mathcal J_{l,-q}=\{r:r\neq q,\ c_r=l\}$.

\paragraph{Difference from the unhierarchical model.}
The sufficient statistics in this section are computed from the auxiliary
positions $r_q$ rather than the LSIRM positions $z_q$.  Because $r_q$ is a
latent variable updated each sweep (see \eqref{eq:r-update}), the NIW
posterior shifts respond to where the conditional distribution of $r_q$
places mass, not directly to the LSIRM-fitted $z_q$.

\section{NIW posterior predictive allocation update for $c_q$}

Marginalizing $(\mu_l,\Sigma_l)$ for label allocation, the posterior
predictive density of $r_q$ given the other members of cluster $l$ is the
Student-$t$ density
\begin{equation}
    r_q\mid r_{\mathcal J_{l,-q}}
    \sim
    t_{\nu_{l,-q}-d+1}
    \!\left(
    m_{l,-q},
    \frac{\kappa_{l,-q}+1}{\kappa_{l,-q}(\nu_{l,-q}-d+1)}S_{l,-q}
    \right).
    \label{eq:niw-predictive}
\end{equation}
The collapsed single-site label update is
\begin{equation}
    P(c_q=l\mid c_{-q},r)
    \propto
    (n_{l,-q}+e_0)\,
    t_{\nu_{l,-q}-d+1}
    \!\left(
    r_q;
    m_{l,-q},
    \frac{\kappa_{l,-q}+1}{\kappa_{l,-q}(\nu_{l,-q}-d+1)}S_{l,-q}
    \right),
    \label{eq:c-update}
\end{equation}
for $l=1,\dots,K^\star$.  This is a Gibbs update on $r$, so the acceptance
probability is one after sampling from the normalized probabilities in
\eqref{eq:c-update}.

\paragraph{Why $r$ and not $z$.}
The collapsed allocation is performed on the auxiliary position $r$ because
the cluster component is defined as $r_q\mid c_q=l\sim N_d(\mu_l,\Sigma_l)$.
Performing the allocation on $z_q$ would require integrating $\Sigma_b$
into the predictive, which is not per-cluster (it is global) and would
break the partition factorization needed by split--merge moves.

\section{Posterior draw of mixture component parameters}

After updating the labels, draw component parameters from their NIW
posterior on $r$:
\begin{align}
    \Sigma_l\mid r,c
    &\sim \IW(\nu_l,S_l),
    \label{eq:sigma-posterior}\\
    \mu_l\mid \Sigma_l,r,c
    &\sim N_d\!\left(m_l,\frac{1}{\kappa_l}\Sigma_l\right),
    \qquad l=1,\dots,K^\star.
    \label{eq:mu-posterior}
\end{align}
For empty components ($n_l=0$), this reduces to a draw from the NIW prior.

After all $\Sigma_l$ have been refreshed, sample $S_0$ from its Wishart
full conditional:
\begin{equation}
    S_0\mid\{\Sigma_l\}_{l=1}^{K^\star}
    \sim
    W\!\left(
    \nu_{S_0}+K^\star\nu_0,\;
    \bigl(\Lambda_{S_0}^{-1}+\textstyle\sum_{l=1}^{K^\star}\Sigma_l^{-1}\bigr)^{-1}
    \right).
    \label{eq:S0-gibbs}
\end{equation}
The cached log-determinant $\log|S_0|$ used in the NIW marginal likelihood
\eqref{eq:niw-marginal} for split--merge moves must be refreshed after every
$S_0$ draw.

The mixture weights may be sampled for monitoring or storage:
\begin{equation}
    \rho\mid c\sim\Dir(e_0+n_1,\dots,e_0+n_{K^\star}).
    \label{eq:rho-posterior}
\end{equation}

\section{Updates for the auxiliary position $r_q$ and the anchor covariance $\Sigma_b$}

\subsection{Per-item conjugate update of $r_q$}

Conditional on $z_q$, $\Sigma_b$, $c_q=l$, $\mu_l$, $\Sigma_l$, the auxiliary
position $r_q$ has Gaussian likelihood from \eqref{eq:z-given-r} and Gaussian
prior from \eqref{eq:r-given-cluster}.  The full conditional is
$N_d(\widehat r_q,\Lambda_q^{-1})$:
\begin{equation}
    \Lambda_q
    =
    \Sigma_b^{-1}+\Sigma_{l}^{-1},
    \qquad
    \widehat r_q
    =
    \Lambda_q^{-1}\bigl(\Sigma_b^{-1}z_q+\Sigma_{l}^{-1}\mu_l\bigr),
    \qquad l=c_q.
    \label{eq:r-update}
\end{equation}
This is a per-item independent Gibbs update; all $P$ updates can be performed
in any order.

\subsection{Conjugate update of $\Sigma_b$ (full IW form)}

Pooling the $z_q-r_q$ deviations across all items,
\begin{equation}
    W_b=\sum_{q=1}^{P}(z_q-r_q)(z_q-r_q)\trans,
    \label{eq:Wb}
\end{equation}
the full conditional of $\Sigma_b$ under the IW prior \eqref{eq:sigma-b-prior}
is
\begin{equation}
    \Sigma_b\mid z,r,c\sim\IW(\nu_b+P,\;S_b+W_b).
    \label{eq:sigma-b-update}
\end{equation}
The degrees of freedom $\nu_b+P$ scales with the total item count, so
$\Sigma_b$ is informed by all clusters jointly.

\paragraph{Wishart hyperprior on $S_b$.}
If \eqref{eq:Sb-hyperprior} is used, the conjugate full conditional is
\begin{equation}
    S_b\mid\Sigma_b
    \sim
    W\!\left(\nu_{S_b}+\nu_b,\;
    (\Lambda_{S_b}^{-1}+\Sigma_b^{-1})^{-1}\right).
    \label{eq:Sb-update}
\end{equation}

\subsection{Scalar isotropic alternative ($\Sigma_b=\sigma_b^2 I_d$)}

Under the recommended default of Section~\ref{sec:identifiability}, the
$z_q\mid r_q$ likelihood becomes $N_d(z_q;r_q,\sigma_b^2 I_d)$ and
\eqref{eq:r-update}, \eqref{eq:sigma-b-update} simplify to
\begin{align}
    \Lambda_q&=\sigma_b^{-2}I_d+\Sigma_{l}^{-1},
    \qquad
    \widehat r_q=\Lambda_q^{-1}\bigl(\sigma_b^{-2}z_q+\Sigma_{l}^{-1}\mu_l\bigr),\\
    \sigma_b^2\mid z,r&\sim\IG\!\left(a_b+\tfrac{Pd}{2},\;
    b_b+\tfrac{1}{2}\sum_{q=1}^{P}\norm{z_q-r_q}_2^2\right).
    \label{eq:sigma-b-scalar-update}
\end{align}
The collapsed label update \eqref{eq:c-update} and split--merge moves are
unchanged.

\section{Collapsed partition target for split--merge moves}

The collapsed partition target on $r$ has the same form as in the
unhierarchical model:
\begin{equation}
    \pi(c\mid r)
    \propto
    \prod_{l=1}^{K^\star}
    \frac{\Gamma(n_l+e_0)}{\Gamma(e_0)}
    m(r_{\mathcal J_l}),
    \label{eq:collapsed-partition-target}
\end{equation}
where $m(r_A)=\int\prod_{q\in A}N_d(r_q;\mu,\Sigma)p(\mu,\Sigma)\,d\mu\,d\Sigma$
is the NIW marginal likelihood evaluated on the auxiliary positions:
\begin{align}
    \log m(r_A)
    &=-\frac{n_Ad}{2}\log\pi
    +\frac{d}{2}\{\log\kappa_0-\log\kappa_A\}
    +\frac{\nu_0}{2}\log|S_0|
    -\frac{\nu_A}{2}\log|S_A|
    \notag\\
    &\qquad+\log\Gamma_d\!\left(\frac{\nu_A}{2}\right)
    -\log\Gamma_d\!\left(\frac{\nu_0}{2}\right),
    \label{eq:niw-marginal}
\end{align}
with $\log\Gamma_d(a)=\frac{d(d-1)}{4}\log\pi+\sum_{h=1}^d\log\Gamma\!\left(a+\frac{1-h}{2}\right)$.
Therefore
\begin{equation}
    \log\pi(c\mid r)
    =
    \sum_{l=1}^{K^\star}
    \!\left[\log\Gamma(n_l+e_0)-\log\Gamma(e_0)
    +\log m(r_{\mathcal J_l})\right]+\mathrm{const}.
    \label{eq:log-collapsed-partition-target}
\end{equation}
Split--merge moves are performed on $c$ conditional on the current $r$.
The global covariance $\Sigma_b$ is not integrated out (it is a separate
Gibbs step), so the partition target factorizes cleanly across clusters.

\section{Metropolis--Hastings updates for LSIRM parameters}

The respondent-effect, threshold, and respondent-position updates are
unchanged from the unhierarchical model.  The only change is in the
item-position update of $b_j^{(\ell)}=z_q$.

\subsection{Respondent effect $\alpha_i^{(\ell)}$}

Propose $\alpha_i^{(\ell)\prime}=\alpha_i^{(\ell)}+\varepsilon_\alpha$,
$\varepsilon_\alpha\sim N(0,s_\alpha^2)$, with log acceptance ratio
\begin{equation}
    \log R_\alpha
    =
    \sum_{j:(i,j)\in\mathcal O_\ell}
    \!\left[\log f_\ell\{y_{ij}^{(\ell)}\!\mid\!\eta_{ij}^{(\ell)\prime},\xi_\ell\}
    -\log f_\ell\{y_{ij}^{(\ell)}\!\mid\!\eta_{ij}^{(\ell)},\xi_\ell\}\right]
    -\frac{(\alpha_i^{(\ell)\prime})^2-(\alpha_i^{(\ell)})^2}{2\sigma_{\alpha,\ell}^2}.
    \label{eq:alpha-mh-ratio}
\end{equation}

\subsection{Item difficulty $\beta_j^{(\ell)}$}

Non-ordinal layers: propose
$\beta_j^{(\ell)\prime}=\beta_j^{(\ell)}+\varepsilon_\beta$,
$\varepsilon_\beta\sim N(0,s_\beta^2)$, with
\begin{equation}
    \log R_\beta
    =
    \sum_{i:(i,j)\in\mathcal O_\ell}
    \!\left[\log f_\ell\{\cdot'\}-\log f_\ell\{\cdot\}\right]
    -\frac{(\beta_j^{(\ell)\prime})^2-(\beta_j^{(\ell)})^2}{2\sigma_{\beta,\ell}^2}.
    \label{eq:beta-mh-ratio}
\end{equation}

Ordinal layers: propose each threshold by Gaussian random walk, reject any
move violating the ordering, and otherwise use
\begin{equation}
    \log R_{\beta,k}
    =
    \sum_{i:(i,j)\in\mathcal O_\ell}\!\left[\log f_\ell\{\cdot'\}-\log f_\ell\{\cdot\}\right]
    -\frac{(\beta_{j,k}^{(\ell)\prime})^2-(\beta_{j,k}^{(\ell)})^2}{2\sigma_{\beta,\ell}^2}.
    \label{eq:beta-thr-mh-ratio}
\end{equation}

\subsection{Respondent latent position $a_i$}

Propose $a_i'=a_i+\varepsilon_a$, $\varepsilon_a\sim N_d(0,s_a^2I_d)$, and
aggregate likelihood contributions across all layers:
\begin{equation}
    \log R_a
    =
    \sum_{\ell=1}^L\sum_{j:(i,j)\in\mathcal O_\ell}\!\left[\log f_\ell\{\cdot'\}-\log f_\ell\{\cdot\}\right]
    -\tfrac{1}{2}\bigl\{\norm{a_i'}_2^2-\norm{a_i}_2^2\bigr\}.
    \label{eq:a-mh-ratio}
\end{equation}

\subsection{Item latent position $b_j^{(\ell)}=z_q$ (modified)}

Let $q=q(\ell,j)$ and let $r_q$ be the current auxiliary position.  Propose
\begin{equation}
    z_q'=z_q+\varepsilon_z,\qquad \varepsilon_z\sim N_d(0,s_z^2 I_d).
\end{equation}
The log acceptance ratio combines the LSIRM likelihood contribution with the
\emph{anchor} contribution from \eqref{eq:z-given-r}:
\begin{equation}
    \log R_z
    =
    \sum_{i:(i,j)\in\mathcal O_\ell}
    \!\left[\log f_\ell\{y_{ij}^{(\ell)}\mid \eta_{ij}^{(\ell)\prime},\xi_\ell\}
    -\log f_\ell\{y_{ij}^{(\ell)}\mid \eta_{ij}^{(\ell)},\xi_\ell\}\right]
    +\log N_d(z_q';r_q,\Sigma_b)-\log N_d(z_q;r_q,\Sigma_b),
    \label{eq:z-mh-ratio}
\end{equation}
where $\eta_{ij}^{(\ell)\prime}$ is the layer-$\ell$ linear predictor
evaluated at $z_q'$.

\paragraph{Comparison with the unhierarchical $b$ update.}
In the unhierarchical model, the anchor was $\bigl(\mu_{c_q},\Sigma_{c_q}\bigr)$.
Here the anchor is $(r_q,\Sigma_b)$: the target shifts from the cluster mean
to the item-specific latent $r_q$, and the precision is the global pooled
$\Sigma_b^{-1}$.  Cluster identity affects the anchor only \emph{indirectly}
through the conditional posterior of $r_q$ in \eqref{eq:r-update}, which
itself blends $z_q$ and $\mu_{c_q}$ in a Gaussian-conjugate way.  This
indirection is the mechanism that decouples cluster tightness from anchor
strength.

\section{Conjugate updates for variance parameters}

The respondent-effect variance update is unchanged:
\begin{equation}
    \sigma_{\alpha,\ell}^2\mid \alpha^{(\ell)}
    \sim\IG\!\left(A_{\alpha,\ell}+\tfrac{n}{2},\;
    B_{\alpha,\ell}+\tfrac{1}{2}\sum_{i=1}^{n}(\alpha_i^{(\ell)})^2\right),
    \qquad \ell=1,\dots,L.
    \label{eq:sigma-alpha-update}
\end{equation}
The cluster-side scale $S_0$ is updated via \eqref{eq:S0-gibbs}.  The
anchor covariance $\Sigma_b$ is updated via \eqref{eq:sigma-b-update} (full
form) or \eqref{eq:sigma-b-scalar-update} (isotropic form).  The
hyperparameter $S_b$ (full form only) is updated via \eqref{eq:Sb-update}.

\section{Jain--Neal restricted Gibbs split--merge moves on $r$}

Split--merge moves are Metropolis--Hastings updates on $c$ conditional on
the current auxiliary positions $r=(r_1,\dots,r_P)$.  The target is the
collapsed partition density in \eqref{eq:log-collapsed-partition-target}.
The structure is identical to the unhierarchical model, with $z\to r$.

\subsection{Anchor selection}

Select two distinct global item indices $u$ and $v$ uniformly from
$\{1,\dots,P\}$.  If $c_u=c_v$, attempt a split move; if $c_u\neq c_v$,
attempt a merge move.  The anchor selection probability is symmetric and
cancels in the Metropolis--Hastings ratio.

\subsection{Split proposal}

Suppose $c_u=c_v=a$ and let $A=\{q:c_q=a\}$.  Choose an empty component
label $b\in\{l:n_l=0\}$ uniformly; if no empty component exists, skip the
split.  Initialize two temporary clusters by forcing $c_u'=a$, $c_v'=b$.
For each $q\in A\setminus\{u,v\}$, assign $q$ to either $a$ or $b$ via a
restricted NIW-collapsed Gibbs scan on $r$:
\begin{align}
    w_q(a)&\propto(n_{a,-q}^{\mathrm{temp}}+e_0)\,p(r_q\mid r_{a,-q}^{\mathrm{temp}}),\\
    w_q(b)&\propto(n_{b,-q}^{\mathrm{temp}}+e_0)\,p(r_q\mid r_{b,-q}^{\mathrm{temp}}),
\end{align}
with the predictive density given by the NIW Student-$t$ in
\eqref{eq:niw-predictive}.  Let $p_q^{\mathrm{fwd}}$ be the normalized
probability of the assignment selected during the forward scan; then
\begin{equation}
    \log q(c'\mid c)
    =-\log E(c)+\sum_{q\in A\setminus\{u,v\}}\log p_q^{\mathrm{fwd}},
    \label{eq:split-forward-probability}
\end{equation}
where $E(c)=|\{l:n_l=0\}|$.  The reverse merge is deterministic, so
$\log q(c\mid c')=0$ and
\begin{equation}
    \alpha_{\mathrm{split}}
    =\min\!\left\{1,\exp[\log\pi(c'\mid r)-\log\pi(c\mid r)-\log q(c'\mid c)]\right\}.
    \label{eq:split-acceptance}
\end{equation}

\subsection{Merge proposal}

Suppose $c_u=a$, $c_v=b$, $a\neq b$, and let $A=\{q:c_q=a\}$,
$B=\{q:c_q=b\}$.  The proposal merges $B$ into $A$:
$c_q'=a$ for $q\in A\cup B$, so $\log q(c'\mid c)=0$.  The reverse split
of $A\cup B$ uses the same anchors $u,v$.  Let $p_q^{\mathrm{rev}}$ be the
normalized probability of assigning $q$ to its original side under the
reverse restricted Gibbs scan; then
\begin{equation}
    \log q(c\mid c')
    =-\log E(c')+\sum_{q\in(A\cup B)\setminus\{u,v\}}\log p_q^{\mathrm{rev}},
    \label{eq:merge-reverse-probability}
\end{equation}
where $E(c')$ is the number of empty components after the merge, and
\begin{equation}
    \alpha_{\mathrm{merge}}
    =\min\!\left\{1,\exp[\log\pi(c'\mid r)-\log\pi(c\mid r)+\log q(c\mid c')]\right\}.
    \label{eq:merge-acceptance}
\end{equation}

\section{Posterior cluster summaries}

Cluster summaries are label-switching invariant.  The primary summary is the
posterior similarity matrix
\begin{equation}
    C_{qr}=P(c_q=c_r\mid Y),
    \qquad q,r=1,\dots,P,
\end{equation}
updated online by $C_{qr}^{\mathrm{raw}}\leftarrow C_{qr}^{\mathrm{raw}}+\ind(c_q=c_r)$
after burn-in, then normalized by the saved-iteration count.  Also monitor
the trace of the occupied-component count $K_+^{(t)}=|\{l:n_l^{(t)}>0\}|$.

\paragraph{Item-level interpretation.}
The label $c_q$ is an item-level cluster assignment that applies to both
$r_q$ and $z_q$.  Visualizations of the LSIRM positions $z_q$ coloured by
$c_q$ are valid; the cluster boundary in $z$-space is fuzzier than the
boundary in $r$-space by an amount controlled by $\Sigma_b$, reflecting the
LSIRM-side anchor noise that the model has explicitly accommodated.

\section{One full MCMC sweep}

\begin{algorithm}[H]
\caption{One MCMC sweep: LSIRM updates and auxiliary positions}
\small
\begin{algorithmic}[1]
\State \textbf{Input:} data $Y$, layer maps $q\mapsto(\ell(q),j(q))$, hyperparameters, proposal scales, $K^\star$, $M_{\mathrm{SM}}$.
\State Update layer-specific nuisance parameters $\xi_\ell$ (e.g.\ $\gamma_\ell$, $\sigma_0^2$, $\lambda_{ij}^{(\ell)}$ for continuous layers; $\{\kappa_j\}$ for count layers), if present.
\For{$i=1,\dots,n$}
    \State Propose $a_i'$ and accept/reject using \eqref{eq:a-mh-ratio} (shared across layers; prior variance fixed at $1$).
\EndFor
\For{$\ell=1,\dots,L$}
    \For{$i=1,\dots,n$}
        \State Propose $\alpha_i^{(\ell)\prime}$ and accept/reject using \eqref{eq:alpha-mh-ratio}.
    \EndFor
    \For{$j=1,\dots,P_\ell$}
        \State Propose $\beta_j^{(\ell)\prime}$ via \eqref{eq:beta-mh-ratio} (non-ordinal), or each threshold via \eqref{eq:beta-thr-mh-ratio} (ordinal, with order check).
        \State Let $q=q(\ell,j)$ and propose $z_q'=b_j^{(\ell)\prime}$.
        \State Accept/reject $z_q'$ using \eqref{eq:z-mh-ratio} with anchor $(r_q,\Sigma_b)$.
    \EndFor
    \State Update $\sigma_{\alpha,\ell}^2$ via \eqref{eq:sigma-alpha-update}.
\EndFor
\State Construct global item positions $z_q=b_{j(q)}^{(\ell(q))}$ for $q=1,\dots,P$.
\For{$q=1,\dots,P$}
    \State Update $r_q\mid z_q,c_q,\mu_{c_q},\Sigma_{c_q},\Sigma_b$ via the conjugate Gaussian draw \eqref{eq:r-update}.
\EndFor
\For{$q=1,\dots,P$}
    \State Remove item $q$ from its current cluster.
    \State Compute NIW leave-one-out predictive weights in \eqref{eq:c-update} using $r_q$ for $l=1,\dots,K^\star$.
    \State Sample $c_q$ from the normalized weights.
\EndFor
\end{algorithmic}
\end{algorithm}

\begin{algorithm}[H]
\caption{One MCMC sweep: split--merge, mixture parameters, anchor covariance, storage}
\small
\begin{algorithmic}[1]
\For{$m=1,\dots,M_{\mathrm{SM}}$}
    \State Select two anchors $u,v$ uniformly at random.
    \If{$c_u=c_v$}
        \State Attempt a restricted-Gibbs split move on $r$ and accept/reject via \eqref{eq:split-acceptance}.
    \Else
        \State Attempt a merge move on $r$ and accept/reject via \eqref{eq:merge-acceptance}.
    \EndIf
\EndFor
\For{$l=1,\dots,K^\star$}
    \State Compute NIW posterior parameters $(\kappa_l,\nu_l,m_l,S_l)$ from \eqref{eq:niw-posterior-params} using $r$.
    \State Draw $\Sigma_l\mid r,c$ via \eqref{eq:sigma-posterior}.
    \State Draw $\mu_l\mid \Sigma_l,r,c$ via \eqref{eq:mu-posterior}.
\EndFor
\State Draw $S_0\mid\{\Sigma_l\}$ via \eqref{eq:S0-gibbs} and refresh $\log|S_0|$ cache.
\State Compute $W_b=\sum_{q=1}^P (z_q-r_q)(z_q-r_q)\trans$ and draw $\Sigma_b$ via \eqref{eq:sigma-b-update} \emph{(full form)} or $\sigma_b^2$ via \eqref{eq:sigma-b-scalar-update} \emph{(isotropic form)}.
\If{Wishart hyperprior on $S_b$ is used (full form only)}
    \State Draw $S_b\mid \Sigma_b$ via \eqref{eq:Sb-update}.
\EndIf
\State Optionally draw $\rho\mid c$ via \eqref{eq:rho-posterior}.
\If{iteration is after burn-in}
    \State Store $c$, $K_+$, $z$, $r$, $\Sigma_b$, and selected LSIRM parameters.
    \State Update posterior similarity counts $C_{qr}^{\mathrm{raw}}\leftarrow C_{qr}^{\mathrm{raw}}+\ind(c_q=c_r)$.
\EndIf
\end{algorithmic}
\end{algorithm}

\section{Implementation notes and recommended defaults}
\label{sec:implementation-notes}

\paragraph{Initial values.}
\begin{itemize}[leftmargin=1.5em]
    \item $z_q^{(0)}$: small Gaussian noise, $z_q^{(0)}\sim N_d(0,0.1\,I_d)$, or initialize from a low-rank embedding of the response matrix.
    \item $r_q^{(0)}=z_q^{(0)}$ (start with no $z$--$r$ split; $\Sigma_b$ will learn the appropriate scale).
    \item $c_q^{(0)}$: random assignment with each $c_q^{(0)}\sim \mathrm{Uniform}\{1,\dots,K_{\mathrm{init}}\}$ for some $K_{\mathrm{init}}\leq K^\star$.
    \item $\Sigma_b^{(0)}=0.05\,I_d$ (or $\sigma_b^{2,(0)}=0.05$).
\end{itemize}

\paragraph{Recommended NIW and IW hyperparameters ($d=2$).}
We recommend default values aligned with the empirically validated
operating point of the unhierarchical model on real and simulated
multilayer data: $\kappa_0=1.0$, $\nu_0=d+10=12$,
$S_0^{(0)}=0.434\,I_d$ (giving
$E[\Sigma_l]=S_0^{(0)}/(\nu_0-d-1)\approx 0.048\,I_d$, comparable to
the within-cluster scale observed in the unhierarchical fit),
$\nu_b=d+10=12$, $S_b^{(0)}=0.09\,I_d$ (giving
$E[\Sigma_b]\approx 0.01\,I_d$), $e_0=0.1$.  Wishart hyperpriors as in
\eqref{eq:S0-hyperprior}, \eqref{eq:Sb-hyperprior} with
$\nu_{S_0}=\nu_{S_b}=d+2$, $\Lambda_{S_0}=S_0^{(0)}/\nu_{S_0}$,
$\Lambda_{S_b}=S_b^{(0)}/\nu_{S_b}$.  Isotropic alternative
\eqref{eq:sigma-b-scalar-prior}: $a_b=3$, $b_b=0.05$, giving
$E[\sigma_b^2]=0.025$ with finite variance.

\paragraph{Why these values.}
A small $\kappa_0$ ($\leq 0.01$) makes the empty-cluster Student-$t$
predictive in \eqref{eq:c-update} highly diffuse (sd
$\propto\sqrt{1/\kappa_0}$ per dim), which inflates the marginal
likelihood of split moves and biases $K_+$ upward; $\kappa_0=1.0$
anchors $\mu_l$ to $m_0$ with prior sd
$\sim\sqrt{S_0/[\kappa_0(\nu_0-d-1)]}\approx 0.22$ per dim, comparable
to the expected within-cluster scale.  A small $\nu_0$ ($=d+2$) gives
$E[\Sigma_l]=S_0/(\nu_0-d-1)$ with infinite variance, leaving
$\Sigma_l$ effectively undetermined for moderate $n_l$; $\nu_0=d+10$
sharpens the prior so that splits producing sub-clusters with
implausibly small $\Sigma_l$ are penalized appropriately.  The
hierarchical extension does not change the role of $\kappa_0$ or the
predictive in \eqref{eq:c-update}, so these considerations carry over
unchanged.

\paragraph{Split--merge frequency.}
The $r_q$ update in \eqref{eq:r-update} pulls $r_q$ toward the cluster
mean of its currently assigned cluster $c_q^{(t-1)}$ before
$c_q^{(t)}$ is resampled.  An incorrectly assigned item therefore
receives a misaligned $r_q$, which biases the subsequent collapsed
allocation \eqref{eq:c-update} toward the same misassignment,
producing a sticky-cluster regime that single-site Gibbs alone cannot
escape.  Empirical experience with the unhierarchical model under
similar scale calibration suggests $M_{\mathrm{SM}}\geq 100$
split--merge attempts per sweep are required for adequate partition
mixing once $b$ moves via random-walk MH (rather than the conjugate
Gibbs available in factor-augmented variants).  The same default is
recommended here, with the per-iter $c_q$-change rate monitored as a
sticky-cluster diagnostic.

\paragraph{Diagnostics.}
\begin{itemize}[leftmargin=1.5em]
    \item Trace of $\tr(\Sigma_b)$ (or $\sigma_b^2$): should stabilize at a moderate non-zero level.  A near-zero $\Sigma_b$ indicates the model is collapsing to a no-anchor regime; a very large $\Sigma_b$ indicates a weakly identified anchor.
    \item Ratio $\tr(\Sigma_b)/\min_l\tr(\Sigma_l)$: a sustained value $\gtrsim 10$ signals that $\Sigma_b$ is absorbing cluster shape rather than acting as global LSIRM-side noise, particularly when small clusters are present (Section~\ref{sec:identifiability}).
    \item Pairwise distance scatter: plot $\norm{z_q-z_r}_2$ against $\norm{r_q-r_r}_2$.  In a healthy fit the scatter should be roughly linear; severe S-curve distortion indicates that more prior mass on larger $\Sigma_b$ may be needed.
    \item Trace of $K_+$: should mix across plausible values; pinned behaviour indicates split--merge moves are too rarely accepted.
    \item Per-iteration $c_q$-change rate: the fraction of items that change cluster between consecutive saved iterations.  Values near zero alongside small split--merge acceptance signal a sticky-cluster trap; consider increasing $M_{\mathrm{SM}}$ or revisiting the $\Sigma_b$ scale.
\end{itemize}

\paragraph{Numerical and post-processing notes.}
\begin{itemize}[leftmargin=1.5em]
    \item In \eqref{eq:r-update}, compute
    $\Lambda_q^{-1}\bigl(\Sigma_b^{-1}z_q+\Sigma_{l}^{-1}\mu_l\bigr)$
    via Cholesky solve rather than explicit matrix inversion; with
    potentially ill-conditioned $\Sigma_l$ for small or empty
    clusters, direct inversion is numerically unstable.
    \item Cache $\Sigma_b^{-1}$ and $\log|\Sigma_b|$ once per sweep
    after the $\Sigma_b$ Gibbs draw; both are used in the $z$-MH ratio
    \eqref{eq:z-mh-ratio} and in the $r_q$ update for every $q$.
    Cache per-cluster $\Sigma_l^{-1}$ similarly after each
    $\Sigma_l$ refresh.
    \item Procrustes alignment for posterior summaries should be
    applied jointly to $\{a_i,b_j^{(\ell)},r_q,\mu_l\}$ under a single
    orthogonal transformation to preserve the hierarchical
    relationship \eqref{eq:z-given-r}--\eqref{eq:r-given-cluster}.
    Aligning $z$ and $r$ separately would distort the $z_q-r_q$
    deviations entering \eqref{eq:Wb} and bias the post-hoc
    $\Sigma_b$ summary.
\end{itemize}

\paragraph{Computational cost per sweep.}
The hierarchical extension adds (i) $P$ Gaussian conjugate draws for $\{r_q\}$, each $O(d^3)$ (Cholesky of $\Lambda_q$, $d=2$); (ii) one IW draw for $\Sigma_b$, $O(d^3)$ after computing $W_b$ in $O(Pd^2)$; (iii) optionally one Wishart draw for $S_b$.  Relative to the unhierarchical sweep these are negligible.  All other steps (NIW posterior, collapsed allocation, split--merge, LSIRM MH) have unchanged per-step cost.

\section{Summary of changes from the unhierarchical model}

\begin{center}
\begin{tabular}{p{0.30\linewidth}p{0.30\linewidth}p{0.30\linewidth}}
\toprule
Component & Unhierarchical & Hierarchical (this document) \\
\midrule
Item position prior & $b_j^{(\ell)}\sim N(\mu_{c_q},\Sigma_{c_q})$ & $z_q\mid r_q\sim N(r_q,\Sigma_b)$, $r_q\mid c_q\sim N(\mu_{c_q},\Sigma_{c_q})$ \\
Anchor in $b$ MH ratio & $(\mu_{c_q},\Sigma_{c_q})$ & $(r_q,\Sigma_b)$ \\
Cluster allocation target & $z_q$ & $r_q$ \\
NIW posterior data & $z$ & $r$ \\
Split--merge target & $\pi(c\mid z)$ & $\pi(c\mid r)$ \\
Anchor strength source & per-cluster $\Sigma_{c_q}$ & global pooled $\Sigma_b$ \\
Cluster shape carrier & $\Sigma_{c_q}$ & $\Sigma_{c_q}$ (unchanged) \\
Marginal $z\mid c=l$ & $N(\mu_l,\Sigma_l)$ & $N(\mu_l,\Sigma_l+\Sigma_b)$ \\
New variables & --- & $\{r_q\}_{q=1}^P$, $\Sigma_b$ (and $S_b$) \\
\bottomrule
\end{tabular}
\end{center}

The cluster-shape heterogeneity carried by the per-cluster $\Sigma_l$ is preserved.  The anchor pull experienced by $z_q$ during its Metropolis--Hastings update is decoupled from cluster identity: it depends on $r_q$ (which adapts via \eqref{eq:r-update}) and the global $\Sigma_b$ (which pools information across all $P$ items via \eqref{eq:sigma-b-update}), neither of which collapses with cluster tightness.  All cluster-allocation machinery is identical in form to the unhierarchical model after replacing $z$ with $r$, so existing implementations can be extended with minimal structural change.

\end{document}